\documentclass[11pt]{article}
\usepackage[T1]{fontenc}
\usepackage{lmodern,amsmath,amssymb,amsthm,mathtools,microtype}
\usepackage[margin=1.15in]{geometry}
\usepackage[hidelinks]{hyperref}
\newtheorem{theorem}{Theorem}[section]
\newtheorem{proposition}[theorem]{Proposition}
\newtheorem{definition}[theorem]{Definition}
\newtheorem{problem}[theorem]{Open Problem}
\newtheorem{example}[theorem]{Example}
\newcommand{\cN}{\mathcal N}
\title{Positive Functionals and Dagger Transfer}
\author{Raeez Lorgat}
\date{}
\hypersetup{
  pdftitle={Positive Functionals and Dagger Transfer},
  pdfauthor={Raeez Lorgat},
  pdfsubject={Pullback of positive forms on possibly nonunital star algebras}
}
\begin{document}
\maketitle

\begin{abstract}
A positive linear functional on a complex $*$-algebra pulls back along a
$*$-homomorphism.  We record the associated sesquilinear form, null space,
and algebraic GNS quotient, and prove that the induced map of GNS spaces is an
isometric injection.  The result is formal: by itself it neither defines a
Weil test algebra nor proves positivity of an explicit-formula functional.
The additional arithmetic realization is therefore stated as a separate,
typed open problem.
\end{abstract}

\section{Positive functionals without a unit}

Throughout, a complex $*$-algebra need not be unital or normed.

\begin{definition}
A complex-linear functional $\omega:\mathcal B\to\mathbb C$ is
\emph{positive} if
\[
 \omega(b^*b)\in\mathbb R_{\geq0}\qquad(b\in\mathcal B).
\]
Its associated sesquilinear form, linear in the first variable, is
\[
 \langle b,c\rangle_\omega=\omega(c^*b),
\]
and its null space is
\[
 \cN_\omega=\{b\in\mathcal B:\omega(b^*b)=0\}.
\]
The algebraic GNS space is the pre-Hilbert space
$\mathcal H_\omega^0=\mathcal B/\cN_\omega$ with the induced form.  Its
Hilbert completion, when desired, is denoted $\mathcal H_\omega$.
\end{definition}

Positivity and polarization imply that the displayed form is Hermitian and
satisfies Cauchy--Schwarz.  In particular, $\cN_\omega$ is a linear subspace.
It is also a left ideal: if $b\in\cN_\omega$ and $a\in\mathcal B$, then
\[
 \omega((ab)^*ab)=\langle a^*ab,b\rangle_\omega=0
\]
by Cauchy--Schwarz.  Thus left multiplication defines the algebraic GNS
representation
\[
 \pi_\omega(a)[b]=[ab]
 \quad\text{on }\mathcal H_\omega^0.
\]
No boundedness assertion is made without a norm and additional hypotheses.

\section{Dagger transfer}

\begin{theorem}[Dagger transfer]
Let $T:\mathcal A\to\mathcal B$ be a $*$-homomorphism of complex
$*$-algebras, let $\omega:\mathcal B\to\mathbb C$ be a positive linear
functional, and put $W=\omega\circ T$.  Then $W$ is positive and
\[
 \langle a,c\rangle_W
   =\langle T(a),T(c)\rangle_\omega,
 \qquad
 \cN_W=T^{-1}(\cN_\omega).                                   \label{pullback}
\]
Consequently
\[
 \widetilde T:\mathcal H_W^0\longrightarrow\mathcal H_\omega^0,
 \qquad [a]\longmapsto[T(a)],                                \label{gnsmap}
\]
is a well-defined isometric injection.  It intertwines the algebraic GNS
actions in the sense that
\[
 \widetilde T\,\pi_W(a)=\pi_\omega(T(a))\,\widetilde T
 \qquad(a\in\mathcal A).
\]
\end{theorem}

\begin{proof}
For $a,c\in\mathcal A$,
\[
 \langle a,c\rangle_W
 =\omega(T(c^*a))
 =\omega(T(c)^*T(a))
 =\langle T(a),T(c)\rangle_\omega.
\]
Taking $c=a$ proves positivity and both equalities in
\eqref{pullback}.  The null-space equality makes \eqref{gnsmap} well defined
and injective, while the form identity makes it isometric.  Finally,
\[
 \widetilde T\bigl(\pi_W(a)[c]\bigr)
 =[T(ac)]=[T(a)T(c)]
 =\pi_\omega(T(a))\widetilde T[c].
\]
\end{proof}

\begin{example}[A finite-dimensional carrier]
Let $\mathcal B=M_n(\mathbb C)$ with its adjoint involution, let
$\rho\in M_n(\mathbb C)$ be positive semidefinite, and define
$\omega_\rho(b)=\operatorname{Tr}(\rho b)$.  Then
\[
 \omega_\rho(b^*b)
   =\operatorname{Tr}\bigl((b\rho^{1/2})^*(b\rho^{1/2})\bigr)\geq0.
\]
For every $*$-homomorphism $T:\mathcal A\to M_n(\mathbb C)$, dagger transfer
gives
\[
 W(a^*a)=\|T(a)\rho^{1/2}\|_{\mathrm{HS}}^2.
\]
This example exhibits all carriers and maps explicitly.
\end{example}

\begin{proposition}[Formal realizability is equivalent to positivity]
For a linear functional $W$ on a complex $*$-algebra $\mathcal A$, the
following are equivalent:
\begin{enumerate}
\item $W$ is positive;
\item there exist a complex $*$-algebra $\mathcal B$ and maps
\[
 \omega:\mathcal B\longrightarrow\mathbb C,
 \qquad T:\mathcal A\longrightarrow\mathcal B,
\]
where $\omega$ is a positive linear functional, $T$ is a
$*$-homomorphism, and $W=\omega\circ T$.
\end{enumerate}
\end{proposition}

\begin{proof}
The theorem proves $(2)\Rightarrow(1)$.  If $W$ is positive, take
$\mathcal B=\mathcal A$, $\omega=W$, and $T=\operatorname{id}_{\mathcal A}$.
\end{proof}

Thus the bare existence of a dagger realization does not furnish an
independent positivity argument; the positivity of its target functional has
to come from additional structure.

\section{The arithmetic realization is additional structure}

For a fixed completed $L$-function, a \emph{normalized Weil datum} will mean
the following explicitly declared objects:
\begin{enumerate}
\item a complex test-function $*$-algebra
      $(\mathcal A_L,\star,*)$;
\item a convergent Hermitian linear functional
      $W_L:\mathcal A_L\to\mathbb C$ written as a specified sum of its polar,
      archimedean, and prime-power local terms;
\item a precisely stated zero-location assertion $\mathsf{RH}_L$ and a proof
      of the equivalence
      \[
      \mathsf{RH}_L
      \quad\Longleftrightarrow\quad
      W_L(f^*\star f)\geq0\quad(f\in\mathcal A_L).
      \]
\end{enumerate}
The test class, involution, convolution, transform convention, local
measures, gamma factors, and any vanishing or moment conditions are part of
the datum.  Weil's original formulations are the primary sources for such
arithmetic positivity criteria for Hecke
$L$-functions~\cite{Weil1952,Weil1972}.  No particular normalization is
selected here.

\begin{problem}[Independent dagger realization of a normalized Weil datum]
Given a normalized Weil datum
$(\mathcal A_L,\star,*,W_L,\mathsf{RH}_L)$, construct a complex $*$-algebra
$\mathcal B$, a positive linear functional $\omega$ on $\mathcal B$, and a
$*$-homomorphism $T:\mathcal A_L\to\mathcal B$ such that
\[
 W_L=\omega\circ T.
\]
Derive $\mathcal B$, $T$, and the positivity of $\omega$ from the local
arithmetic or geometric input without using $\mathsf{RH}_L$, its zero set, or
an assertion equivalent to its Weil positivity criterion.  Prove the
displayed identity term by term with the chosen polar, archimedean, and
prime-power normalizations.
\end{problem}

The theorem would then transfer positivity to $W_L$.  The proposition shows
why the independence condition is essential: without it, the requested
factorization merely restates the desired positivity.

\begin{thebibliography}{9}

\bibitem{Weil1952}
A.~Weil,
\emph{Sur les ``formules explicites'' de la th\'eorie des nombres premiers},
Commun. S\'em. Math. Univ. Lund, Tome suppl\'ementaire (1952), 252--265.

\bibitem{Weil1972}
A.~Weil,
\emph{Sur les formules explicites de la th\'eorie des nombres},
Math. USSR-Izv. \textbf{6} (1972), 1--17.
\href{https://doi.org/10.1070/IM1972v006n01ABEH001866}
{doi:10.1070/IM1972v006n01ABEH001866}.

\end{thebibliography}

\end{document}
